\makeatletter
\newcommand{\DeclareLatinAbbrev}[2]{%
  \DeclareRobustCommand{#1}{%
  \@ifnextchar\cite{\textit{#2.\,}}{%  if next is a \cite then insert a period and a thin space
    \@ifnextchar{.}{\textit{#2}}{%
      \@ifnextchar{,}{\textit{#2.}}{%
        \@ifnextchar{!}{\textit{#2.}}{%
          \@ifnextchar{?}{\textit{#2.}}{%
            \@ifnextchar{)}{\textit{#2.}}{%
              {\textit{#2.,\ }}}}}}}}}%
}
\makeatother
\DeclareLatinAbbrev{\eg}{e.g}
\DeclareLatinAbbrev{\Eg}{E.g}
\DeclareLatinAbbrev{\ie}{i.e}
\DeclareLatinAbbrev{\Ie}{I.e}
\DeclareLatinAbbrev{\etc}{etc}
\DeclareLatinAbbrev{\etal}{et~al}


\def\first {(\textit{i})\xspace}
\def\Second{(\textit{ii})\xspace}
\def\third {(\textit{iii})\xspace}
\def\fourth{(\textit{iv})\xspace}
\def\fifth {(\textit{v})\xspace}
\def\sixth {(\textit{vi})\xspace}
\def\seventh{(\textit{vii})\xspace}
\def\eighth{(\textit{viii})\xspace}


%%% custom definitions
%% Coloring the links!
\newcommand\myshade{75} % Usage: red!\myshade!black

\definecolor{ForestGreen} {RGB}{34,  139,  34}
\definecolor{HeraldRed2}   {rgb}{0.81, 0.12, 0.15}

\newcommand{\refscolor} {blue}
\newcommand{\linkscolor}{HeraldRed2}
\newcommand{\urlscolor} {ForestGreen}

%% Some definitions of used colors
%\definecolor{darkblue}{rgb}{0.0,0.0,0.3} %% define a color called darkblue
%\definecolor{darkred}{rgb}{0.4,0.0,0.0}
%\definecolor{red}{rgb}{0.7,0.0,0.0}
%\definecolor{lightgrey}{rgb}{0.8,0.8,0.8} 
%\definecolor{grey}{rgb}{0.6,0.6,0.6}
%\definecolor{darkgrey}{rgb}{0.4,0.4,0.4}
%\definecolor{aqua}{rgb}{0.0, 1.0, 1.0}

% For runin headings
\newcommand{\smartparagraph}[1]{\vspace{.05in}\noindent\textbf{#1}}

%% Table of Contents (ToC) depth 
\setcounter{secnumdepth}{5} % how many sectioning levels to assign numbers to
\setcounter{tocdepth}{4}    % how many sectioning levels to show in ToC

%% Limit hyphenation
\hyphenpenalty=9000
\tolerance=5000
% Reduce hyphenation as much as possible:
%\hyphenpenalty=15000
%\tolerance=1000

% For notes by the authors to themselves
\newcommand*{\todoinline}[1]{\textcolor{red}{TODO: #1}}

% define a command to set BibTeX and BibLaTeX nicely
\DeclareRobustCommand{\BibTeX}{B\textsc{ib}\!T\kern-.1667em\lower.5ex\hbox{E}\kern-.125emX}

\DeclareRobustCommand{\BibLaTeX}{\textsc{Bib}\LaTeX}

%\DeclareUnicodeCharacter{2003}{\quad}

% Use grouping for numbers that are 4 or more digits long
\sisetup{group-minimum-digits=4}

\definecolor{kth-blue}{RGB/cmyk}{25,84,166/0.849,0.494,0,0.349}
\colorlet{kth-blue80}{kth-blue!80!}
\colorlet{kth-blue40}{kth-blue!40!}

%% Secondary colors
\definecolor{kth-lightblue}{RGB/cmyk}{36,160,216/0.833,0.259,0,0.153}
\colorlet{kth-lightblue80}{kth-lightblue!80!}
\colorlet{kth-lightblue40}{kth-lightblue!40!}

\definecolor{rowtint}{RGB}{230, 237, 246}


\definecolor{pageRefColor}{HTML}{D83389} 
